\documentclass[12pt]{article}

\usepackage{graphicx}
\usepackage{amssymb}
\usepackage{amsmath}
\usepackage[margin=1in]{geometry}
\usepackage{fancyhdr}
\usepackage{enumerate}
\usepackage[shortlabels]{enumitem}
\usepackage{algorithm}
\usepackage{algpseudocode}

\pagestyle{fancy}
\fancyhead[l]{Li Yifeng}
\fancyhead[c]{Homework \#1}
\fancyhead[r]{\today}
\fancyfoot[c]{\thepage}
\renewcommand{\headrulewidth}{0.2pt}
\setlength{\headheight}{15pt}

\newcommand{\td}{\text{DNF}}

\begin{document}
	
	\section*{Exercise 5}
	
	\noindent Define disjunctive normal form formulas. When would it be helpful to have a formula in $\td$? Sketch in pseudocode an algorithm to turn any propositional formula into a $\td$ formula.
	
	\bigskip
	
	\begin{enumerate}[start=1,label={\bfseries Part \arabic*:},leftmargin=0in]
		\bigskip\item Define disjunctive normal form formulas.
		
		\subsection*{Answer}
			
			A formula is in disjunctive normal form ($\td$) if it is a finite disjunction of conjunctions of \emph{literals}, where a literal is either an atomic proposition $p$ or its negation $\lnot p$. Formally:
			\[
			\varphi \text{ is $\td$ } \iff \varphi \equiv \bigvee_{i=1}^k \bigwedge_{j=1}^{m_i} \ell_{ij},
			\]
			with each $\ell_{ij}\in\{p,\lnot p\}$. By convention, $\top$ is the empty conjunction and $\bot$ is the empty disjunction.
			
		\bigskip\item When would it be helpful to have a formula in $\td$?
		
		\subsection*{Answer}
			
			Disjunctive normal form formulas are useful to evaluate satisfiability and represent logical circuits.
			
		\bigskip\item Sketch in pseudocode an algorithm to turn any propositional formula into a $\td$ formula.
			
		\subsection*{Answer}
			
			\begin{algorithm}
				
				\caption{\td(F)}
				
				\begin{algorithmic}
					
					\State $F \gets \Call{remove\_implications}{F}$ \Comment{remove $\to$, $\leftrightarrow$}
					\State $F \gets \Call{push\_negation}{F}$ \Comment{move $\lnot$ inward (De Morgan)}
					\State $F \gets \Call{distribute\_or}{F}$ \Comment{distribute $\lor$ over $\land$}
					\State \Return \Call{simplified}{F}
					
				\end{algorithmic}
				
			\end{algorithm}
			
	\end{enumerate}
	
\end{document}